% Bagian: turing-machines
% Bab: machines-computations
% Subbagian: turing-machines

\documentclass[../../../include/open-logic-section]{subfiles}

\begin{document}

\olfileid[id]{tur}{mac}{tur}
\olsection{Mesin Turing}

\begin{explain}
Definisi formal mengenai apa yang membentuk suatu mesin Turing tampak
abstrak, tetapi sebenarnya sederhana: definisi itu hanya mengemas ke dalam satu
struktur matematis seluruh informasi yang diperlukan untuk menentukan cara
kerja mesin Turing. Informasi ini meliputi (1) keadaan-keadaan yang dapat
ditempati mesin, (2) simbol-simbol yang boleh berada pada pita, (3) keadaan
tempat mesin harus bermula, dan (4) himpunan instruksi mesin.
\end{explain}

\begin{defn}[Mesin Turing]
Sebuah \emph{mesin Turing} $M$ adalah tuple $\langle Q, \Sigma, q_0,
\delta\rangle$ yang terdiri atas
\begin{enumerate}
\item suatu himpunan berhingga \emph{keadaan}~$Q$,
\item suatu \emph{alfabet} berhingga $\Sigma$ yang mencakup $\TMendtape$ dan
  $\TMblank$,
\item suatu \emph{keadaan awal}~$q_0 \in Q$,
\item suatu \emph{himpunan instruksi} berhingga~$\delta\colon Q \times \Sigma
  \pto Q \times \Sigma \times \{\TMleft, \TMright, \TMstay\}$.
\end{enumerate}
Fungsi parsial~$\delta$ juga disebut \emph{fungsi transisi} dari~$M$.
\end{defn}

\begin{explain}
Kita mengasumsikan bahwa pita tak hingga ke satu arah saja. Karena itu,
berguna untuk menetapkan simbol khusus~$\TMendtape$ sebagai penanda ujung kiri
pita. Penanda ini memudahkan program mesin Turing mengetahui kapan mesin
``terancam'' keluar dari pita. Kita dapat mengasumsikan bahwa simbol ini tidak
pernah ditimpa, yakni bahwa $\delta(q,\TMendtape) = \tuple{q', \TMendtape, x}$
jika $\delta(q,\TMendtape)$ terdefinisi. Sebagian buku teks menggunakan
asumsi ini; kita tidak. Anda cukup berhati-hati ketika menyusun mesin Turing
agar mesin itu tidak pernah menimpa~$\TMendtape$. Selain itu, ada keadaan
ketika mengizinkan penimpaan tersebut memberikan keluwesan yang berguna.
\end{explain}

\begin{ex}
\emph{Mesin Genap:} Secara formal, mesin genap adalah kuadruple
$\tuple{Q, \Sigma, q_0, \delta}$ dengan
\begin{align*}
Q & = \{ q_0, q_1 \} \\
\Sigma & = \{ \TMendtape, \TMblank, \TMstroke \}, \\
\delta(q_0, \TMstroke) & = \tuple{q_1, \TMstroke, \TMright},\\
\delta(q_1, \TMstroke) & = \tuple{q_0, \TMstroke, \TMright},\\
\delta(q_1, \TMblank)  & = \tuple{q_1, \TMblank, \TMright}.
\end{align*}
\end{ex}

\end{document}
